\subsection{MINLPLib models}\label{sec:minlplib}

On MINLPLib models the cuts were neutral for SCIP: they changed neither which
models it solved nor its own solving time. Part~3A repeats the 30 models of campaign~2 (ten per stratum, selected by a
hash ordering from 422 eligible models with at most 120 variables) with a
300-second budget, seeds 0, 1 and~2, and root runs. Part~3B uses a sample
selected by structure: from the 392 eligible models not used in campaign~2,
the 111 on which discovery finds a block that the \code{auto} rule admits,
ranked by a hash, of which we take the first 30, with seeds 0 and~1. The root
runs of both parts include mode \code{all-diag}, whose raised limits test
whether the default limits, rather than the cuts, bound their effect. SCIP
solves most of these models in about a second, so they can show validity and
cost but hardly a benefit. Part~4D addresses this. Its pool consists of the
322 MINLPLib models with 121 to 1{,}000 variables that are not convex by the
library's metadata, contain a quadratic or polynomial function and were not
used before. Of these, 85 have a block that \code{auto} admits, and 58 of those
are not solved by SCIP within 60 seconds; Part~4D takes the first 20 of the 58
in a fixed hash order (Appendix~\ref{app:instances}). Table~\ref{tab:minlplib} summarizes the full
runs; Tables~\ref{tab:partA}, \ref{tab:partB} and~\ref{tab:partD} in
Appendix~\ref{app:instances} give the results per model.

\begin{table}[t]
\centering\footnotesize
\caption{Full runs on MINLPLib models (300\,s). Times are shifted geometric
means (shift 1\,s) over the runs solved by every mode of the part; ``SCIP
excl.'' is SCIP's solving time minus the separator callback; ``ratio'' is the
median per-run ratio of the total time to SCIP. ``Bounds'' counts
final dual bounds better/worse than SCIP at relative tolerance
$10^{-4}$ over all paired runs. In Part~4D only two models were solved, and
times are not comparable.}\label{tab:minlplib}
\begin{tabular}{@{}llrrrrrr@{}}
\toprule
Part & Mode & Solved & Cuts (runs) & Time (s) & SCIP excl.\ (s) & Ratio & Bounds\\
\midrule
3A (30 models, 3 seeds) & SCIP & 80/90 & --- & 1.05 & 0.99 & --- & ---\\
  & \code{all}  & 80/90 & 151 (33) & 1.27 & 0.99 & 1.34 & 2/4\\
  & \code{auto} & 80/90 & 37 (15)  & 1.21 & 0.99 & 1.29 & 3/3\\
\midrule
3B (30 models, 2 seeds) & SCIP & 52/60 & --- & 1.33 & 1.27 & --- & ---\\
  & \code{all}  & 52/60 & 179 (40) & 1.87 & 1.25 & 2.66 & 2/0\\
  & \code{auto} & 52/60 & 145 (38) & 1.82 & 1.23 & 2.55 & 2/0\\
\midrule
4D (20 models, 1 seed) & SCIP & 2/20 & --- & & & & ---\\
  & \code{all}  & 2/20 & 14 (2) & & & & 5/1\\
  & \code{auto} & 2/20 & 14 (2) & & & & 4/0\\
  & SCIP-extra & 1/20 & --- & & & & 7/7\\
\bottomrule
\end{tabular}
\end{table}

\paragraph{Parts 3A and 3B.} In every part and seed, all modes solved the same
models. In Part~3A, cuts were added on 11 of the 30 models in mode \code{all}
and on 5 in mode \code{auto}. The final dual bounds of the 90 paired runs were
better in 2 and worse in 4 for \code{all}, and better in 3 and worse in 3 for
\code{auto}; between seeds, the baseline itself differed in 3 of 30
comparisons, so these differences are within run-to-run variation. In
Part~3B, cuts were added on 20 of the 30 models; the final bounds were better
in 2 of 60 comparisons in each cut mode, both on \code{kall\_circles\_c8a} at
the time limit, and never worse. In modes \code{all} and \code{auto}, the cuts changed
the root bound visibly on four models of Part~3B and on none of Part~3A. They closed the root gap of
\code{pooling\_bental4tp} in every cut mode and that of
\code{pooling\_bental4pq} in modes \code{auto} and \code{all-diag}, which were
then solved at the root, and they improved the root bounds of \code{ex3\_1\_4}
(from $-6$ to $-5.79$, and to $-5.69$ with raised limits; the optimum is $-4$)
and of \code{pointpack04} (a maximization model, for which a smaller bound is
stronger). In two root runs, \code{auto} ended with a worse root
bound, on \code{pointpack04} and on \code{pooling\_haverly2pq}, where the bound
fell from $-617.7$ to $-857.1$ (the optimum is $-600$). Valid cuts cannot
weaken an LP relaxation, but they change SCIP's later separation and cut
selection, and the bound at the end of the root node need not improve.
Raising the limits at the root (\code{all-diag}) added more cuts, 403 in
Part~3A and 358 in Part~3B, but improved the root bound beyond the tolerance on
only one more model in each part than mode \code{all} (\code{cvxnonsep\_psig20r}
in Part~3A and \code{pooling\_bental4pq} in Part~3B). In the post hoc root runs of Part~3P, whole-row
directions changed the number of cuts in 10 of 180 cut-mode runs and no root
bound. None of these effects changed a full solve. The cut
modes were 15--20\% slower than SCIP in shifted geometric mean on
Part~3A and 37--41\% on Part~3B, and slower per run by median factors of 1.3 and
2.5--2.7; for comparison, baseline runs of Part~3A that differ only in SCIP's
seed differed per run by a median factor of 1.03 (range 0.36 to 2.38 over all
pairs of seeds). SCIP's
own solving time, without the separator callback, was unchanged (Table~\ref{tab:minlplib}), and the node counts summed over the
commonly solved runs differed by less than 3\%. The cuts are therefore neutral
for SCIP on these models, and the slowdown is the cost of the separator
(Section~\ref{sec:funnel}).

\paragraph{Presolve and the stored-row check (Part 4B2).} In campaign~3 the
stored-row check discarded 20\% (\code{all}) and 36\% (\code{all-diag}) of the
certified and violated rows of the Part~3B root runs. Reruns that recorded the
cause showed that presolve had aggregated or fixed a block variable in 74--85\%
of these cases, so that SCIP stored the cut in other columns. Part~4B2 repeats
the root runs with SCIP's aggregation and multi-aggregation disabled in every
mode. The rejected share fell to 5\% and 9.5\%; the remaining rejections come
from variables that presolve fixed and from SCIP's own coefficient handling
(integer snapping and coefficients below $10^{-9}$ dropped). Against SCIP with
the same presolve setting, mode \code{all} improved the root bound on 3 models
and worsened it on none, and \code{all-diag} improved it on 6 models
(including \code{pooling\_bental4tp}, \code{pooling\_bental4pq},
\code{ex3\_1\_4}, \code{sep1} and \code{pointpack04}) and worsened it on one
(\code{pooling\_haverly1tp}). Disabling aggregation, however, weakened SCIP's
own root bound on 6 of the 30 models, so the setting is not a free remedy. On
\code{kall\_circlespolygons\_c1p12}, where 127 of 152 violated rows had been
rejected before, 104 cuts now reached the LP, and the root bound stayed at $0$
against an optimum of $0.34$. Trying the support of each whole row first
(\code{all-diag-rowdir}) gave the same cuts and bounds on every model.

\paragraph{SCIP's disabled separators.} SCIP-extra, with SCIP's
disabled-by-default nonconvex separators switched on at the root, improved the
root bound on 12 of the 30 Part~3B models and worsened it on one
(\code{sep1}). It closed the root gap completely on five pooling and small
nonconvex models and by more than 95\% on two more, and these seven include
the four models on which the certified cuts had an effect; it solved 13 models
at the root against 8. Intersection cuts and interminor cuts were productive
on 26 and 21 models, edge-concave cuts on none. On these models, SCIP's own
alternatives are clearly stronger than the certified cuts.

\paragraph{Larger models (Part 4D).} In mode \code{all} the separator stopped
during block discovery on 12 of the 20 models in the root runs (11 in the full
runs), exactly those on which
discovery had taken more than the one-second allowance in the scan (1.0 to
7.7\,s). On six of the other eight it completed discovery and added no cut:
every certified support either fell below the violation threshold or, on three
models, was rejected by the stored-row check because presolve had fixed or
aggregated a block variable (11 of the 25 violated rows). Discovery also
found more blocks than the cap of 32 on 6 of these 8 models, and on 10 of the
20 with the cap of 128 in mode \code{all-diag}; the blocks tried were then
those with the smallest variable indices. With raised limits
at the root (\code{all-diag}), the separator added 891 cuts on 13 models and
improved the root bound on one model (\code{multiplants\_mtg1c}, closing 4\%
of its root gap), while the root bound was worse on five, by at most 3\% of the
gap; the variant with whole-row directions behaved the same. In the full runs
the cut modes solved the same two models as the baseline. All final-bound
differences in Table~\ref{tab:minlplib} occurred on models that received no
cut, so they reflect the changed timing of a time-limited run, not the cuts.
SCIP-extra improved the root bound on 8 of the 20 models, worsened it on 3,
and in full runs lost one of the two solved models. On models of this size the
separator with default limits stops during discovery on most models, and
where it gets past discovery its cuts are not stronger than SCIP's
relaxation.
